\underline{Exercise $4.4$}

\begin{enumerate}
\item[5.] (e)  \[ \begin{aligned}
 & A=\left[\begin{array}{c}
-2 \\
4 \\
7
\end{array}\right]_{3 \times 1} \quad B=\left[\begin{array}{lll}
3 & 6 & -2
\end{array}\right]_{1 \times 3} \\~\\
& C=A B=\left[\begin{array}{ccc}
-6 & -12 & 4 \\
12 & 24 & -8 \\
21 & 42 & -14
\end{array}\right]_{3 \times 3} \\~\\
& D=B A=[3 \times-2+6 \times 4+-2 \times 7]=[4]_{1 \times 1}
\end{aligned} \]

% Question 7
\item[7.] In example 5, $
\begin{aligned} x=\left[\begin{array}{l}x_{1} \\x_{2}\end{array}\right] \text{ and }A=\left[\begin{array}{cc}a_{11} & 0 \\0 & a_{22}\end{array}\right] \end{aligned} $.
In which case,
\[
\begin{aligned}
& x^{\prime} A x=\left[\begin{array}{ll}x_{1} & x_{2}\end{array}\right]\left[\begin{array}{cc}a_{11} & 0 \\0 & a_{22}\end{array}\right]\left[\begin{array}{l}x_{1} \\x_{2}\end{array}\right]_{2 \times 1} \\
& =\left[\begin{array}{lll}a_{11} x_{1} & a_{22} x_{2}\end{array}\right]_{1 \times 2}\left[\begin{array}{l}x_{1} \\x_{2}\end{array}\right]_{2 \times 1} \\
& =\underbrace{a_{11} x_{1}^{2}+a_{22} x_{2}^{2}}_{\text{Weighted sum of squares}}=\sum_{i=1}^{2} a_{i i} x_{i}^{2}
\end{aligned}
\]

So $x^{\prime} A x$ represents a weighted sum of squares where $a_{11}, a_{22}$ are weights.\\

But now what if $A=\left[\begin{array}{ll}a_{11} & a_{12} \\ a_{21} & a_{22}\end{array}\right]$. In this case,

\[
x^{\prime} A x=\left[\begin{array}{ll}
x_{1} & x_{2}
\end{array}\right]\left[\begin{array}{ll}
a_{11} & a_{12} \\
a_{21} & a_{22}
\end{array}\right]_{2 \times 2}\left[\begin{array}{l}
x_{1} \\
x_{2}
\end{array}\right]_{2 \times 1}
\]

\[
\begin{aligned}
& =\left[\begin{array}{rr}
a_{11} x_{1}+a_{21} x_{2} & a_{12} x_{1}+a_{22} x_{2}
\end{array}\right]\left[\begin{array}{l}
x_{1} \\
x_{2}
\end{array}\right]_{2 \times 1} \\
& =a_{11} x_{1}^{2}+a_{21} x_{1} x_{2}+a_{12} x_{1} x_{2}+a_{22} x_{2}^{2} \\
& =a_{11} x_{1}^{2}+\left(a_{21}+a_{12}\right) x_{1} x_{2}+a_{22} x_{2}^{2}
\end{aligned}
\]

So $x^{\prime} A x$ no longer represents a weighted sum of squares. \\

You can check that the associative law i.e.
\[
\left(x^{\prime} A\right) x=x^{\prime}(A x)
\]
will apply in both cases (after all, its a law!) as all products are possible.
\end{enumerate}

\newpage
%%%%% Exercise 4.5

\underline{Exercise $4.5$}
\begin{enumerate}
% Question 1
\item[1.] \[A=\left[\begin{array}{ccc}-1 & 5 & 7 \\0 & -2 & 4\end{array}\right]_{2 \times 3} \quad \quad b=\left[\begin{array}{c}9 \\6 \\0\end{array}\right]_{3 \times 1} \quad \quad x=\left[\begin{array}{l}x_{1} \\x_{2}\end{array}\right]_{2\times 1}\]

\begin{enumerate}

\item $ \begin{aligned} A I=\left[\begin{array}{ccc}-1 & 5 & 7 \\0 & -2 & 4\end{array}\right]\left[\begin{array}{lll}1 & 0 & 0 \\0 & 1 & 0 \\0 & 0 & 1\end{array}\right]_{3 \times 3}  =\left[\begin{array}{ccc}-1 & 5 & 7 \\0 & -2 & 4\end{array}\right]=A \end{aligned}$  \\~\\

\item $ \begin{aligned} IA =\left[\begin{array}{ll}
1 & 0 \\
0 & 1
\end{array}\right]_{2\times 2}\left[\begin{array}{ccc}
-1 & 5 & 7 \\
0 & -2 & 4
\end{array}\right]  =\left[\begin{array}{rrr}
-1 & 5 & 7 \\
0 & -2 & 4
\end{array}\right] \end{aligned}$ \\~\\

\item $ \begin{aligned} Ix=\left[\begin{array}{ll}1 & 0 \\ 0 & 1\end{array}\right]_{2 \times 2}\left[\begin{array}{l}x_{1} \\ x_{2}\end{array}\right]=\left[\begin{array}{l}x_{1} \\ x_{2}\end{array}\right]\end{aligned}$ \\~\\

\item $ \begin{aligned} x^{\prime} I=\left[\begin{array}{ll}x_{1} & x_{2}\end{array}\right]\left[\begin{array}{ll}1 & 0 \\ 0 & 1\end{array}\right]_{2 \times 2}=\left[\begin{array}{ll}x_{1} & x_{2}\end{array}\right]\end{aligned}$ \\~\\
\end{enumerate}

% Question 4
\item[4.] Let's start with a $2 \times 2$ diagonal matrix

\[
\left[\begin{array}{cc}
a_{11} & 0 \\
0 & a_{22}
\end{array}\right]\left[\begin{array}{cc}
a_{11} & 0 \\
0 & a_{22}
\end{array}\right]=\left[\begin{array}{cc}
a_{11}^{2} & 0 \\
0 & a_{22}^{2}
\end{array}\right]
\]

$x=x^{2}$ for only $x=0,1$ so $a_{11}$ and $a_{22}$ can either be 0 or 1 . So we can have the following $2 \times 2$ idempotent diagonal matrices:

\[
\left[\begin{array}{ll}
1 & 0 \\
0 & 0
\end{array}\right],\left[\begin{array}{ll}
1 & 0 \\
0 & 1
\end{array}\right],\left[\begin{array}{ll}
0 & 0 \\
0 & 0
\end{array}\right],\left[\begin{array}{ll}
0 & 0 \\
0 & 1
\end{array}\right]
\]

More generally, for $n \times n$ matrix, there can be $2^{n}$ such matrices. This is because there will be $n$ elements, each of which can take two values.
\end{enumerate}

%%%%% Exercise 4.6

\underline{Exercise $4.6$}
\begin{enumerate}
% Question 2
\item[2.] $
A=\left[\begin{array}{cc}
0 & 4 \\
-1 & 3
\end{array}\right] \quad B=\left[\begin{array}{cc}
3 & -8 \\
0 & 1
\end{array}\right] \quad C=\left[\begin{array}{lll}
1 & 0 & 9 \\
6 & 1 & 1
\end{array}\right]
$

\begin{enumerate}

\item[(a)] \( A+B=\left[\begin{array}{cc}3 & -4 \\ -1 & 4\end{array}\right] \) \\~\\
\( A^{\prime}+B^{\prime}=\left[\begin{array}{rr}0 & -1 \\ 4 & 3\end{array}\right]+\left[\begin{array}{cc}3 & 0 \\ -8 & 1\end{array}\right]=\left[\begin{array}{cc}3 & -1 \\ -4 & 4\end{array}\right] \) \\~\\
So, \((A+B)^{\prime}=A^{\prime}+B^{\prime}\) \\

\item[(b)] \(
A C=\left[\begin{array}{cc}
0 & 4 \\
-1 & 3
\end{array}\right]_{2 \times 2}\left[\begin{array}{lll}
1 & 0 & 9 \\
6 & 1 & 1
\end{array}\right]_{2 \times 3}=\left[\begin{array}{ccc}
24 & 4 & 4 \\
17 & 3 & -6
\end{array}\right]_{2 \times 3}
\)
\vspace{1em}
\(
C^{\prime} A^{\prime}=\left[\begin{array}{ll}
1 & 6 \\
0 & 1 \\
9 & 1
\end{array}\right]_{3 \times 2}\left[\begin{array}{cc}
0 & -1 \\
4 & 3
\end{array}\right]_{2 \times 2}=\left[\begin{array}{cc}
24 & 17 \\
4 & 3 \\
4 & -6
\end{array}\right]_{3 \times 2}
\) \vspace{1em}
So, $(AC)^{\prime}=C^{\prime} A^{\prime}$.
\end{enumerate}
\end{enumerate}

%%%%% Exercise 5.1

\underline{Exercise 5.1}
\begin{enumerate}

% Question 3
\item[3.]
\begin{multicols}{2}
\begin{enumerate}
\item[(a)] Yes
\item[(b)] Yes
\item[(c)] Yes
\item[(d)] No; second row $= -2 \times$ first row.
\end{enumerate}
\end{multicols}

% Question 4
\item[4.] Yes, we get the same answer.

\end{enumerate}
